\documentclass[12pt]{article}

\input{../../../_assets/preambulo.tex}

\begin{document}

    \input{../../../_assets/portada}
    \portadaExamenFotoDif{../../dueA4.jpg}{GKI\\Exam I}{GKI. Exam I}{MidnightBlue}{-9}{30}{2026}{Arturo Olivares Martos}

    \begin{description}
        \item[Subject] Fundamentals of Artificial Intelligence.
        \item[Academic Year] Winter Semester 2025/2026.
        % \item[Degree] ---.
        % \item[Group] ---.
        \item[Lecturer] Gérald Kämmerer.
        \item[Description] Final Exam.
        \item[Date] 05.03.2026.
        \item[Duration] 90 minutes.
    \end{description}
    \newpage

    \section*{Part 1: Short Questions (47 pts.)}

    \begin{ejercicio}[Supervised Learning (3 pts.)]
        What is supervised learning? Mark whether each of the following methods belongs to supervised learning or not.

        \begin{table}[h!]
        \centering
        \begin{tabular}{lcc}
        \hline
        \textbf{Method} & \textbf{Right} & \textbf{Wrong} \\ \hline
        kNN & $\square$ & $\square$ \\
        ID3 & $\square$ & $\square$ \\
        Hill Climbing & $\square$ & $\square$ \\
        OPTICS & $\square$ & $\square$ \\
        CNN & $\square$ & $\square$ \\
        LSTM & $\square$ & $\square$ \\ \hline
        \end{tabular}
        \caption{Classification of Supervised Learning Methods}
        \end{table}
    \end{ejercicio}

    \begin{ejercicio}[PPDAC Cycle (5 pts.)]
        List the steps of the PPDAC cycle and give an application example.
    \end{ejercicio}

    \begin{ejercicio}[Data Aggregation Errors (2 pts.)]
        Name two possible errors that can occur when data is aggregated.
    \end{ejercicio}

    \begin{ejercicio}[Basic Statistics (5 pts.)]
        Given the numbers $A = \{10, 0, 6, 0, 9\}$, compute the median, mean, mode, and variance.
    \end{ejercicio}

    \begin{ejercicio}[Data Encoding Issues (2 pts.)]
        Name two problems that can occur when we encode data.
    \end{ejercicio}

    \begin{ejercicio}[Neural Network Optimization (2 pts.)]
        Name two methods used for optimizing neural networks.
    \end{ejercicio}

    \begin{ejercicio}[Search Algorithms (1 pt.)]
        For which search method is Depth-Limited Search a modification?
    \end{ejercicio}

    \begin{ejercicio}[Genetic Algorithms - Evaluation (2 pts.)]
        Given three chromosomes of length $8$ with the bases $\{a, b, c\}$:
        \begin{itemize}
            \item Chromosome 1: $cabbacba$
            \item Chromosome 2: $aabaabca$
            \item Chromosome 3: $babbaaba$
        \end{itemize}
        Evaluate the chromosomes based on the frequency of $a$ (the more $a$'s, the better). What would the best chromosome from the given set look like, and what does the theoretically best chromosome look like?
    \end{ejercicio}

    \begin{ejercicio}[Genetic Algorithms - Recombination \& Mutation (3 pts.)]
        Given two chromosomes: Chromosome 1 ($acbcaaba$) and Chromosome 2 ($cabaabca$). Apply order-based recombination (OX2) using the sequence $1,2,4,6$. Afterwards, apply a mutation.
    \end{ejercicio}

    \begin{ejercicio}[Uniform Cost Search (6 pts.)]
        Explain Uniform Cost Search using the graph in Figure~\ref{fig:UCS}. Go through the different steps. Note down the path/steps, the respective costs, and the total cost.

        \begin{figure}[h!]
        \centering
        \begin{tikzpicture}[
            node/.style={circle, draw, fill=black, text=white, minimum size=1.2cm, align=center, font=\small\sffamily\bfseries},
            edge/.style={-Stealth, thick}
        ]
        \node[node] (A) at (0,0) {Start A};
        \node[node] (B) at (3,2) {B};
        \node[node] (C) at (3,0) {C};
        \node[node] (D) at (3,-2) {D};
        \node[node] (E) at (6,1.5) {Goal E};
        \node[node] (F) at (6,-1.5) {F};

        \draw[edge, -Stealth, Stealth-] (A) -- node[above left, text=black] {2} (B);
        \draw[edge] (A) -- node[above, text=black] {8} (C);
        \draw[edge] (A) -- node[below left, text=black] {1} (D);
        \draw[edge, -Stealth, Stealth-] (B) -- node[left, text=black] {1} (C);
        \draw[edge] (B) -- node[above, text=black] {0} (E);
        \draw[edge] (B) -- node[above right, text=black] {2} (F);
        \draw[edge, -Stealth, Stealth-] (C) -- node[left, text=black] {3} (D);
        \draw[edge, -Stealth, Stealth-] (C) -- node[above, text=black] {3} (F);
        \end{tikzpicture}
        \caption{Uniform Cost Graph}\label{fig:UCS}
        \end{figure}
    \end{ejercicio}

    \begin{ejercicio}[Regularization Models (5 pts.)]
        List the definitions and differences between Ridge, Lasso, and ElasticNet regularization.
    \end{ejercicio}

    \begin{ejercicio}[Specificity (1 pt.)]
        How is the specificity (TNR) defined, and in which value range does it lie?
    \end{ejercicio}

    \begin{ejercicio}[Confusion Matrix (2 pts.)]
        Write down the general confusion matrix and name its elements.
    \end{ejercicio}

    \begin{ejercicio}[Hierarchical Clustering (2 pts.)]
        Define Single Linkage and Average Linkage in hierarchical clustering, and state the mutual advantages and disadvantages of each.
    \end{ejercicio}

    \begin{ejercicio}[Autoencoders (1 pt.)]
        Name the main components of an autoencoder.
    \end{ejercicio}

    \begin{ejercicio}[NLP Lemmatization (2 pts.)]
        Lemmatize the words in the sentence: \textit{``I am well prepared''}.
    \end{ejercicio}

    \begin{ejercicio}[Neural Network Architecture (3 pts.)]
        Draw a feed-forward neural network with:
        \begin{itemize}
            \item 1 input layer ($1$ neuron)
            \item 2 hidden layers ($2$ neurons each)
            \item 1 output layer ($1$ neuron)
        \end{itemize}
        The first hidden layer must be fully connected to the second hidden layer. You can choose the other connections freely. Additionally, draw and label exactly one direct feedback connection and one indirect feedback connection.
    \end{ejercicio}

    \newpage

    \section*{Part 2: k-Means (13 pts.)}

    \begin{ejercicio}[k-Means Algorithm Steps (3 pts.)]
        Explain the steps involved in performing the k-Means algorithm.
    \end{ejercicio}

    \begin{ejercicio}[Execution of k-Means (10 pts.)]
        The following dataset (Table 2) is provided. We want to cluster these points and find centroids.

        \begin{table}[h!]
        \centering
        \begin{tabular}{|c|c|c|}
        \hline
        \textbf{Point} & \textbf{x} & \textbf{y} \\ \hline
        P1 & 1 & 2 \\ \hline
        P2 & 1 & 4 \\ \hline
        P3 & 2 & 1 \\ \hline
        P4 & 5 & 4 \\ \hline
        P5 & 6 & 5 \\ \hline
        P6 & 6 & 7 \\ \hline
        \end{tabular}
        \caption{Data for k-Means Clustering}
        \end{table}

        Execute the k-Means algorithm step-by-step. Start with $k = 2$ and initial centroids at points P3 and P5. Fully execute the first two iterations. Distances are calculated using Manhattan distance:
        \[ d(P1, P2) = |x_1 - x_2| + |y_1 - y_2| \]

        \textbf{Example calculation for distance between points:}
        \[ d(P1, P2) = |1 - 1| + |2 - 4| = 0 + 2 = 2 \]
    \end{ejercicio}

    \newpage

    \section*{Part 3: Kahn's Algorithm, Entropy \& Random Forest (22 pts.)}

    \begin{ejercicio}[Kahn's Algorithm (6 pts.)]
        Name and explain the steps of Kahn's Algorithm.
    \end{ejercicio}

    \begin{ejercicio}[Feature Types (2 pts.)]
        Consider Table~\ref{tab:customer_dataset}. Which features are categorical and which are nominal?

        \begin{table}[h!]
        \centering
        \begin{tabular}{|l|l|l|l|}
        \hline
        \textbf{Client} & \textbf{Age} & \textbf{Ad} & \textbf{Buy} \\ \hline
        1 & young & yes & yes \\ \hline
        2 & young & no & no \\ \hline
        3 & medium & yes & yes \\ \hline
        4 & old & yes & yes \\ \hline
        5 & old & yes & no \\ \hline
        6 & medium & yes & yes \\ \hline
        \end{tabular}
        \caption{Customer Dataset}
        \label{tab:customer_dataset}
        \end{table}
    \end{ejercicio}

    \begin{ejercicio}[Entropy Calculation (9 pts.)]
        Which feature should you split on first in Table~\ref{tab:customer_dataset} if the column \textit{Buy} is the target? Calculate the entropy for each feature using the formula:
        \[ S(x) = -\sum_{i=1}^{N} p(x_i) \cdot \log_2(p(x_i)) \]
    \end{ejercicio}

    \begin{ejercicio}[Information Gain Range (2 pts.)]
        What is the range of values for Information Gain? Specify its minimum and maximum values.
    \end{ejercicio}

    \begin{ejercicio}[Random Forest Classification (3 pts.)]
        Classify the point $(F_1 = 1, F_2 = 4, F_3 = 7)$ using the Random Forest in Figure~\ref{fig:random_forest}. Provide the explicit result for each tree as well as the final aggregated result. Draw the explicit decision path. Note: the left branch of a node represents \textbf{true}, and the right branch represents \textbf{false}.

        \begin{figure}[h!]
            \centering
            \resizebox{\textwidth}{!}{%
            \begin{tikzpicture}[
                node/.style={circle, draw=black, minimum size=1.2cm, inner sep=0pt, font=\scriptsize\bfseries, align=center},
                leaf/.style={circle, draw=black, minimum size=0.85cm, inner sep=0pt, font=\scriptsize\bfseries},
                edge from parent/.style={draw, -Stealth, ultra thick},
                tree label/.style={font=\small\bfseries, yshift=0.35cm}
            ]

            % ==================== TREE 1 ====================
            \begin{scope}[shift={(0,0)}]
                \node[tree label] at (0,0) {1};
                \node[node] at (0,-0.6) {$F_2 \leq 4.5$}
                    [sibling distance=3.2cm, level distance=1.4cm]
                    child { node[node] {$F_1 \leq 1$}
                        [sibling distance=1.6cm, level distance=1.4cm]
                        child { node[node] {$F_3 > 6$}
                            [sibling distance=1.2cm, level distance=1.4cm]
                            child { node[leaf] {C1} }
                            child { node[leaf] {C3} }
                        }
                        child { node[leaf] {C4} }
                    }
                    child { node[node] {$F_1 > 1$}
                        [sibling distance=1.6cm, level distance=1.4cm]
                        child { node[node] {$F_1 < -1$}
                            [sibling distance=1.2cm, level distance=1.4cm]
                            child { node[leaf] {C3} }
                            child { node[leaf] {C2} }
                        }
                        child { node[leaf] {C1} }
                    };
            \end{scope}

            % ==================== TREE 2 ====================
            \begin{scope}[shift={(6.2,0)}]
                \node[tree label] at (0,0) {2};
                \node[node] at (0,-0.6) {$F_2 > 2$}
                    [sibling distance=2.6cm, level distance=1.4cm]
                    child { node[node] {$F_1 < 0$}
                        [sibling distance=1.5cm, level distance=1.4cm]
                        child { node[node] {$F_3 < 9$}
                            [sibling distance=1.2cm, level distance=1.4cm]
                            child { node[leaf] {C1} }
                            child { node[leaf] {C2} }
                        }
                        child { node[leaf] {C1} }
                    }
                    child { node[node] {$F_2 \leq 1$}
                        [sibling distance=1.2cm, level distance=1.4cm]
                        child { node[leaf] {C4} }
                        child { node[leaf] {C2} }
                    };
            \end{scope}

            % ==================== TREE 3 ====================
            \begin{scope}[shift={(11.8,0)}]
                \node[tree label] at (0,0) {3};
                \node[node] at (0,-0.6) {$F_3 \leq 10$}
                    [sibling distance=1.8cm, level distance=1.4cm]
                    child { node[node] {$F_3 \leq 5$}
                        [sibling distance=1.3cm, level distance=1.4cm]
                        child { node[leaf] {C3} }
                        child { node[leaf] {C1} }
                    }
                    child { node[leaf] {C2} };
            \end{scope}

            % ==================== TREE 4 ====================
            \begin{scope}[shift={(7.2,-5.2)}]
                \node[tree label] at (0,0) {4};
                \node[node] at (0,-0.6) {$F_1 \geq 2$}
                    [sibling distance=1.4cm, level distance=1.4cm]
                    child { node[leaf] {C2} }
                    child { node[leaf] {C5} };
            \end{scope}

            % ==================== TREE 5 ====================
            \begin{scope}[shift={(10.8,-5.2)}]
                \node[tree label] at (0,0) {5};
                \node[node] at (0,-0.6) {$F_2 \geq 4$}
                    [sibling distance=1.4cm, level distance=1.4cm]
                    child { node[leaf] {C4} }
                    child { node[leaf] {C1} };
            \end{scope}

            \end{tikzpicture}
            }
            \caption{Random Forest Classifier}
            \label{fig:random_forest}
        \end{figure}
    \end{ejercicio}

    \newpage

    \section*{Part 4: NLP \& Forward Propagation (8 pts.)}

    \begin{ejercicio}[Bag of Words Encoding (4 pts.)]
        Consider the sentence \textit{``Nobody makes it better than me, nobody.''} and the vocabulary
        $$\{A, be, better, easy, is, it, make, makes, me, nobody, than, then\}$$
        
        Encode the sentence using frequency-based Bag of Words.
    \end{ejercicio}

    \begin{ejercicio}[Forward Propagation Computation (4 pts.)]
        Given the small neural network below, write the calculated output values (after applying the activation functions) into the corresponding boxes. The weights $w$ are provided for all connections.

        \begin{figure}[h!]
            \centering
            \resizebox{\textwidth}{!}{%
            \begin{tikzpicture}[
                neuron/.style={circle, draw, fill=black, text=white, minimum size=1.3cm, font=\small\sffamily\bfseries, align=center},
                box/.style={rectangle, draw, minimum size=0.65cm, thick, fill=white},
                edge/.style={-Stealth, thick},
                weight/.style={font=\small, fill=white, inner sep=1.5pt, rounded corners=1pt, sloped}
            ]

            % --- Nodes Placement ---
            % Input
            \node[font=\large\bfseries] (in) at (-1.5,0) {7};

            % Layer 1: Heaviside
            \node[neuron] (H) at (1,0) {Heaviside};
            \node[box] (H_box) at (2.4,0) {};

            % Layer 2: ReLUs
            \node[neuron] (R1) at (6, 2.8) {ReLU};
            \node[box] (R1_box) at (6, 3.9) {};

            \node[neuron] (R2) at (6, 0) {ReLU};
            \node[box] (R2_box) at (6, 1.1) {};

            \node[neuron] (R3) at (6, -2.8) {ReLU};
            \node[box] (R3_box) at (6, -3.9) {};

            % Layer 3: Sign
            \node[neuron] (S1) at (11.5, 1.6) {Sign};
            \node[box] (S1_box) at (11.5, 2.7) {};

            \node[neuron] (S2) at (11.5, -1.6) {Sign};
            \node[box] (S2_box) at (11.5, -2.7) {};

            % Layer 4: Output ReLU
            \node[neuron] (Out) at (16, 0) {ReLU};
            \node[box] (Out_box) at (17.4, 0) {};


            % --- Connections & Weights ---

            % Input -> Heaviside
            \draw[edge] (in) -- (H);

            % Heaviside Box -> ReLUs
            \draw[edge] (H_box) -- node[weight, above] {$w_2 = 1$} (R1);
            \draw[edge] (H_box) -- node[weight, above] {$w_1 = -2$} (R2);
            \draw[edge] (H_box) -- node[weight, below] {$w_3 = 1$} (R3);

            % ReLUs -> Sign Layer (Weights placed clearly near start or near end, sloped along the line)
            \draw[edge] (R1) -- node[weight, pos=0.2, above] {$w_1 = 2$} (S1);
            \draw[edge] (R1) -- node[weight, pos=0.2, below] {$w_2 = 2$} (S2);

            \draw[edge] (R2) -- node[weight, pos=0.7, above] {$w_3 = 2$} (S1);
            \draw[edge] (R2) -- node[weight, pos=0.7, below] {$w_4 = 2$} (S2);

            \draw[edge] (R3) -- node[weight, pos=0.2, above] {$w_5 = 2$} (S1);
            \draw[edge] (R3) -- node[weight, pos=0.2, below] {$w_6 = 2$} (S2);

            % Sign Layer -> Output
            \draw[edge] (S1) -- node[weight, above] {$w_1 = -2$} (Out);
            \draw[edge] (S2) -- node[weight, below] {$w_2 = 2$} (Out);

            \end{tikzpicture}
            }
            \caption{Forward Propagation Network Schema}
        \end{figure}
    \end{ejercicio}


    % SOLUCION
    \setcounter{section}{0}
    \setcounter{ejercicio}{0}
    
    \newpage\newpage
    \section*{Part 1: Short Questions (47 pts.)}

    \begin{ejercicio}[Supervised Learning (3 pts.)]
        What is supervised learning? Mark whether each of the following methods belongs to supervised learning or not.

        \begin{table}[h!]
        \centering
        \begin{tabular}{lcc}
        \hline
        \textbf{Method} & \textbf{Right} & \textbf{Wrong} \\ \hline
        kNN & $\boxtimes$ & $\square$ \\
        ID3 & $\boxtimes$ & $\square$ \\
        Hill Climbing & $\square$ & $\boxtimes$ \\
        OPTICS & $\square$ & $\boxtimes$ \\
        CNN & $\boxtimes$ & $\square$ \\
        LSTM & $\boxtimes$ & $\square$ \\ \hline
        \end{tabular}
        \caption{Classification of Supervised Learning Methods}
        \end{table}
    \end{ejercicio}

    \begin{ejercicio}[PPDAC Cycle (5 pts.)]
        List the steps of the PPDAC cycle and give an application example.
    \end{ejercicio}

    \begin{ejercicio}[Data Aggregation Errors (2 pts.)]
        Name two possible errors that can occur when data is aggregated.

        When data is aggregated, two possible errors that can occur are:
        \begin{enumerate}
            \item Loss of granularity: Aggregating data can lead to the loss of detailed information, which may obscure important patterns or trends present in the original dataset.
            \item Introduction of bias: Aggregation can introduce bias if the aggregation method does not accurately represent the underlying data distribution, leading to misleading conclusions or interpretations.
        \end{enumerate}
    \end{ejercicio}

    \begin{ejercicio}[Basic Statistics (5 pts.)]
        Given the numbers $A = \{10, 0, 6, 0, 9\}$, compute the median, mean, mode, and variance.
        \begin{itemize}
            \item \ul{Median}: To find the median, we first sort the numbers: $0, 0, 6, 9, 10$. The median is the middle value, which is $6$.
            \item \ul{Mean}: The mean is calculated by summing all the numbers and dividing by the count of numbers. 
            \[ \text{Mean} = \frac{10 + 0 + 6 + 0 + 9}{5} = \frac{25}{5} = 5 \]
            \item \ul{Mode}: The mode is the number that appears most frequently. In this case, the mode is $0$ since it appears twice.
            \item \ul{Variance}: Variance is calculated using the formula:
            \begin{align*}
                \text{Variance} &= \frac{\sum (x_i - \text{mean})^2}{N} = \dfrac{(10-5)^2 + (0-5)^2 + (6-5)^2 + (0-5)^2 + (9-5)^2}{5} =\\&= \dfrac{25 + 25 + 1 + 25 + 16}{5} = \dfrac{92}{5} = 18.4
            \end{align*}
        \end{itemize}
    \end{ejercicio}

    \begin{ejercicio}[Data Encoding Issues (2 pts.)]
        Name two problems that can occur when we encode data.
        \begin{enumerate}
            \item Collision: When encoding data, different inputs may produce the same encoded output, leading to ambiguity and loss of information.
            \item Collinearity: Encoding can introduce multicollinearity, where encoded features are highly correlated, which can negatively impact the performance of certain machine learning models. The hot-encoding method, for example, can lead to collinearity if not handled properly.
        \end{enumerate}
    \end{ejercicio}

    \begin{ejercicio}[Neural Network Optimization (2 pts.)]
        Name two methods used for optimizing neural networks.

        Two common methods for optimizing neural networks are:
        \begin{enumerate}
            \item Gradient Descent: An iterative optimization algorithm used to minimize the loss function by updating the weights in the direction of the negative gradient.
            \item Stochastic Gradient Descent (SGD): A variant of gradient descent that updates the weights using a randomly selected subset of the training data, which can lead to faster convergence and better generalization.
        \end{enumerate}
    \end{ejercicio}

    \begin{ejercicio}[Search Algorithms (1 pt.)]
        For which search method is Depth-Limited Search a modification?

        Depth-Limited Search is a modification of the Depth-First Search (DFS) algorithm. It introduces a depth limit to the search process, preventing the algorithm from exploring nodes beyond a specified depth in the search tree. This modification helps to avoid infinite loops in cases where the search space is infinite or very large.
    \end{ejercicio}

    \begin{ejercicio}[Genetic Algorithms - Evaluation (2 pts.)]
        Given three chromosomes of length $8$ with the bases $\{a, b, c\}$:
        \begin{itemize}
            \item Chromosome 1: $cabbacba$
            \item Chromosome 2: $aabaabca$
            \item Chromosome 3: $babbaaba$
        \end{itemize}
        Evaluate the chromosomes based on the frequency of $a$ (the more $a$'s, the better). What would the best chromosome from the given set look like, and what does the theoretically best chromosome look like?

        \begin{itemize}
            \item Chromosome 1: $cabbacba$ - Frequency of $a$: 3
            \item Chromosome 2: $aabaabca$ - Frequency of $a$: 5
            \item Chromosome 3: $babbaaba$ - Frequency of $a$: 4
        \end{itemize}

        The best chromosome from the given set is Chromosome 2 ($aabaabca$) with a frequency of $5$ $a$'s. The theoretically best chromosome would be one that consists entirely of $a$'s, such as $aaaaaaaa$, which would have a frequency of $8$ $a$'s.
    \end{ejercicio}

    \begin{ejercicio}[Genetic Algorithms - Recombination \& Mutation (3 pts.)]
        Given two chromosomes: Chromosome 1 ($acbcaaba$) and Chromosome 2 ($cabaabca$). Apply order-based recombination (OX2) using the sequence $1,2,4,6$. Afterwards, apply a mutation.
        \begin{table}[ht]
            \centering
            \begin{tabular}{r|cccccccc}
                Parent 1 & a & c & b & c & a & a & b & a \\ \hline
                Parent 2 & c & a & b & a & a & b & c & a \\ \hline
                Offspring & c & c & b & a & a & b & b & a\\ \hline
                Swap Mutation & c & c & b & a & a & b & \textbf{a} & \textbf{b}
            \end{tabular}
            \caption{\centering Example of Order-Based Crossover (OX2), where the fixed positions are 1, 2, 4, and 6.}
            \label{tab:ox2_example}
        \end{table}
        
        This exercise has been solved in Table~\ref{tab:ox2_example}.
    \end{ejercicio}

    \begin{ejercicio}[Uniform Cost Search (6 pts.)]
        Explain Uniform Cost Search using the graph in Figure~\ref{fig:UCS2}. Go through the different steps. Note down the path/steps, the respective costs, and the total cost.

        \begin{figure}[h!]
            \centering
            \begin{tikzpicture}[
                node/.style={circle, draw, fill=black, text=white, minimum size=1.2cm, align=center, font=\small\sffamily\bfseries},
                edge/.style={-Stealth, thick}
            ]
            \node[node] (A) at (0,0) {Start A};
            \node[node] (B) at (3,2) {B};
            \node[node] (C) at (3,0) {C};
            \node[node] (D) at (3,-2) {D};
            \node[node] (E) at (6,1.5) {Goal E};
            \node[node] (F) at (6,-1.5) {F};

            \draw[edge, -Stealth, Stealth-] (A) -- node[above left, text=black] {2} (B);
            \draw[edge] (A) -- node[above, text=black] {8} (C);
            \draw[edge] (A) -- node[below left, text=black] {1} (D);
            \draw[edge, -Stealth, Stealth-] (B) -- node[left, text=black] {1} (C);
            \draw[edge] (B) -- node[above, text=black] {0} (E);
            \draw[edge] (B) -- node[above right, text=black] {2} (F);
            \draw[edge, -Stealth, Stealth-] (C) -- node[left, text=black] {3} (D);
            \draw[edge, -Stealth, Stealth-] (C) -- node[above, text=black] {3} (F);
            \end{tikzpicture}
            \caption{Uniform Cost Graph}\label{fig:UCS2}
        \end{figure}

        We use a priority queue to keep track of the nodes to explore, sorted by their cumulative cost from the start node. The steps are as follows:
        \begin{enumerate}
            \item $PQ = [(A, 0)]$
            \item Expand A: $PQ = [(D, 1), (C, 8)]$
            \item Expand D: $PQ = [(C, 4)]$
            \item Expand C: $PQ = [(B, 5)]$
            \item Expand B: $PQ = [(E, 5), (F, 7)]$
            \item Expand E: Goal reached with total cost = 5.
        \end{enumerate}
    \end{ejercicio}

    \begin{ejercicio}[Regularization Models (5 pts.)]
        List the definitions and differences between Ridge, Lasso, and ElasticNet regularization.

        The normal loss function for a linear regression model is given by:
        \[ L(\theta) = \sum_{i=1}^{N} (y_i - \hat{y}_i)^2 \]

        However, in regularization, we add a penalty term to the loss function to prevent overfitting. The definitions are as follows:
        \begin{itemize}
            \item Lasso (L1 regularization): Given $\lambda\in \bb{R}^+$, the loss function is modified to:
            \[ L(\theta) = \sum_{i=1}^{N} (y_i - \hat{y}_i)^2 + \lambda \sum_{j=1}^{p} |\theta_j| \]

            This penalty encourages sparsity in the model parameters, effectively performing feature selection by driving some coefficients to zero.

            \item Ridge (L2 regularization): Given $\lambda\in \bb{R}^+$, the loss function is modified to:
            \[ L(\theta) = \sum_{i=1}^{N} (y_i - \hat{y}_i)^2 + \lambda \sum_{j=1}^{p} \theta_j^2 \]

            This penalty discourages large coefficients but does not enforce sparsity, meaning all features are retained in the model.

            \item ElasticNet: Given $\lambda_1, \lambda_2\in \bb{R}^+$, the loss function is modified to:
            \[ L(\theta) = \sum_{i=1}^{N} (y_i - \hat{y}_i)^2 + \lambda_1 \sum_{j=1}^{p} |\theta_j| + \lambda_2 \sum_{j=1}^{p} \theta_j^2 \]

            This combines both L1 and L2 penalties, allowing for a balance between sparsity and coefficient shrinkage, making it suitable for situations where there are correlated features.
        \end{itemize}
    \end{ejercicio}

    \begin{ejercicio}[Specificity (1 pt.)]
        How is the specificity (TNR) defined, and in which value range does it lie?
    \end{ejercicio}

    \begin{ejercicio}[Confusion Matrix (2 pts.)]
        Write down the general confusion matrix and name its elements.

        The general confusion matrix is a table that summarizes the performance of a classification model by comparing the predicted labels with the true labels. The coefficient $a_{ij}$ represents the number of instances where the true class is $i$ and the predicted class is $j$. In a binary classification scenario, the confusion matrix is typically represented as follows:
        \[
        \begin{bmatrix}
        \text{True Positive (TP)} & \text{False Negative (FN)} \\
        \text{False Positive (FP)} & \text{True Negative (TN)}
        \end{bmatrix}
        \]
    \end{ejercicio}

    \begin{ejercicio}[Hierarchical Clustering (2 pts.)]
        Define Single Linkage and Average Linkage in hierarchical clustering, and state the mutual advantages and disadvantages of each.

        Both are methods for measuring the distance between clusters in hierarchical clustering:
        \begin{itemize}
            \item Single Linkage: The distance between two clusters is defined as the minimum distance between any single point in one cluster and any single point in the other cluster.
            \begin{equation*}
                d_{\text{single}}(C_1, C_2) = \min_{x \in C_1, y \in C_2} d(x, y)
            \end{equation*}
            This method can lead to the ``chaining effect,'' where clusters may be formed by linking together points that are close to each other, potentially resulting in elongated and less compact clusters.

            \item Average Linkage: The distance between two clusters is defined as the average distance between all pairs of points in the two clusters.
            \begin{equation*}
                d_{\text{average}}(C_1, C_2) = \frac{1}{|C_1||C_2|} \sum_{x \in C_1} \sum_{y \in C_2} d(x, y)
            \end{equation*}
            This method tends to produce more balanced and compact clusters, but it can be computationally more intensive, especially for large datasets, as it requires calculating the average distance between all pairs of points in the clusters.
        \end{itemize}
    \end{ejercicio}

    \begin{ejercicio}[Autoencoders (1 pt.)]
        Name the main components of an autoencoder.

        The main components of an autoencoder are:
        \begin{itemize}
            \item Encoder: The encoder is a neural network that compresses the input data into a lower-dimensional representation, often referred to as the latent space or code. It learns to capture the essential features of the input data while reducing its dimensionality.
            \item Latent Space: The latent space is the compressed representation of the input data produced by the encoder. It contains the most relevant information about the input data and serves as the input for the decoder.
            \item Decoder: The decoder is a neural network that reconstructs the original input data from the latent space representation. It learns to map the compressed features back to the original data space, aiming to minimize the reconstruction error between the input and the output.
        \end{itemize}
    \end{ejercicio}

    \begin{ejercicio}[NLP Lemmatization (2 pts.)]
        Lemmatize the words in the sentence: \textit{``I am well prepared''}.

        The lemmatized form of the words in the sentence \textit{``I am well prepared''} is:
        \begin{itemize}
            \item I $\rightarrow$ I
            \item am $\rightarrow$ be
            \item well $\rightarrow$ well
            \item prepared $\rightarrow$ prepare
        \end{itemize}
    \end{ejercicio}

    \begin{ejercicio}[Neural Network Architecture (3 pts.)]
        Draw a feed-forward neural network with:
        \begin{itemize}
            \item 1 input layer ($1$ neuron)
            \item 2 hidden layers ($2$ neurons each)
            \item 1 output layer ($1$ neuron)
        \end{itemize}
        The first hidden layer must be fully connected to the second hidden layer. You can choose the other connections freely. Additionally, draw and label exactly one direct feedback connection and one indirect feedback connection.\\

        The feed-forward neuronal network is represented in Figure~\ref{fig:ffnn}.
        \begin{figure}
            \centering
            \begin{tikzpicture}[
                    neuron/.style={circle, draw, fill=black, text=white, minimum size=1cm, font=\small\sffamily\bfseries},
                    edge/.style={-Stealth, thick},
                    feedback/.style={edge, dashed, red, thin}, % Estilo para feedback: punteado, rojo
                    label_fb/.style={font=\scriptsize, text=red, fill=white, inner sep=1pt, rounded corners=1pt}
                ]

                % --- Posicionamiento de Nodos ---
                % Capa Entrada
                \node[neuron] (I1) at (0,0) {$I_1$};

                % Capa Oculta 1
                \node[neuron, yshift=1cm] (H11) at (3,0) {$H_{1,1}$};
                \node[neuron, yshift=-1cm] (H12) at (3,0) {$H_{1,2}$};

                % Capa Oculta 2
                \node[neuron, yshift=1cm] (H21) at (6,0) {$H_{2,1}$};
                \node[neuron, yshift=-1cm] (H22) at (6,0) {$H_{2,2}$};

                % Capa Salida
                \node[neuron] (O1) at (9,0) {$O_1$};


                % --- Conexiones Feed-Forward (Estándar, Negras) ---
                % Entrada -> Oculta 1 (Conexiones elegidas libremente)
                \draw[edge] (I1) -- (H11);
                \draw[edge] (I1) -- (H12);

                % Oculta 1 -> Oculta 2 (FULLY CONNECTED)
                \draw[edge] (H11) -- (H21);
                \draw[edge] (H11) -- (H22);
                \draw[edge] (H12) -- (H21);
                \draw[edge] (H12) -- (H22);

                % Oculta 2 -> Salida (Conexiones elegidas libremente)
                \draw[edge] (H21) -- (O1);
                \draw[edge] (H22) -- (O1);


                % --- Conexiones de Feedback (Retroalimentación, Rojas, Punteadas) ---
                % 1. Direct Feedback Connection: O1 -> O1
                \draw[feedback] (O1.north) .. controls +(0.5, 1) and +(-0.5, 1) .. (O1.north)
                    node[label_fb, pos=0.5, anchor=south, label={above:Direct Feedback}] {};

                % 2. Indirect Feedback Connection: H2,2 -> H1,1
                \draw[feedback] (H22.south) .. controls +(0, -1.5) and +(0, -1.5) .. (H12.south)
                    node[label_fb, pos=0.5, anchor=north] {Indirect Feedback};

            \end{tikzpicture}
            \caption{Feed-Forward Neural Network with Feedback Connections.}
            \label{fig:ffnn}
        \end{figure}
    \end{ejercicio}

    \newpage

    \section*{Part 2: k-Means (13 pts.)}

    \begin{ejercicio}[k-Means Algorithm Steps (3 pts.)]
        Explain the steps involved in performing the k-Means algorithm.

        The k-Means algorithm is an iterative clustering algorithm that partitions a dataset into $k$ distinct clusters based on feature similarity. The steps involved in performing the k-Means algorithm are as follows:
        \begin{enumerate}
            \item Initialization: Randomly select $k$ initial centroids from the dataset. These centroids can be chosen randomly or using specific initialization methods like k-means++ to improve convergence.
            \item Assignment Step: Assign each data point to the nearest centroid based on a distance metric (e.g., Euclidean or Manhattan distance). This step forms $k$ clusters, where each cluster contains the data points closest to its corresponding centroid.
            \item Update Step: Recalculate the centroids of each cluster by taking the mean (or median) of all the data points assigned to that cluster. The new centroids represent the updated positions of the clusters.
            \item Convergence Check: Repeat the assignment and update steps until convergence is achieved. Convergence can be defined as no change in cluster assignments or centroids, or when a maximum number of iterations is reached.
        \end{enumerate}
    \end{ejercicio}

    \begin{ejercicio}[Execution of k-Means (10 pts.)]
        The following dataset (Table 2) is provided. We want to cluster these points and find centroids.

        \begin{table}[h!]
        \centering
        \begin{tabular}{|c|c|c|}
        \hline
        \textbf{Point} & \textbf{x} & \textbf{y} \\ \hline
        P1 & 1 & 2 \\ \hline
        P2 & 1 & 4 \\ \hline
        P3 & 2 & 1 \\ \hline
        P4 & 5 & 4 \\ \hline
        P5 & 6 & 5 \\ \hline
        P6 & 6 & 7 \\ \hline
        \end{tabular}
        \caption{Data for k-Means Clustering}
        \end{table}

        Execute the k-Means algorithm step-by-step. Start with $k = 2$ and initial centroids at points P3 and P5. Fully execute the first two iterations. Distances are calculated using Manhattan distance:
        \[ d(P1, P2) = |x_1 - x_2| + |y_1 - y_2| \]

        \textbf{Example calculation for distance between points:}
        \[ d(P1, P2) = |1 - 1| + |2 - 4| = 0 + 2 = 2 \]

        \begin{enumerate}
            \item The assigment step in the first iteration is described in Table~\ref{tab:kmeans_iteration1}.
            \begin{table}[h!]
            \centering
            \begin{tabular}{|c|c|c|c|}
            \hline
            \textbf{Point} & \textbf{Distance to P3} & \textbf{Distance to P5} & \textbf{Assigned Cluster} \\ \hline
            P1 & 2 & 8 & Cluster 1 \\ \hline
            P2 & 4 & 6 & Cluster 1 \\ \hline
            P3 & 0 & 8 & Cluster 1 \\ \hline
            P4 & 6 & 2 & Cluster 2 \\ \hline
            P5 & 8 & 0 & Cluster 2 \\ \hline
            P6 & 10 & 2 & Cluster 2 \\ \hline
            \end{tabular}
            \caption{Assignment Step - Iteration 1}
            \label{tab:kmeans_iteration1}
            \end{table}

            \item The new centroids after the first iteration are calculated as follows:
            \begin{itemize}
                \item Cluster 1 (P1, P2, P3):
                \begin{equation*}
                    \text{New centroid 1} = \left(\frac{1 + 1 + 2}{3}, \frac{2 + 4 + 1}{3}\right) = \left(\frac{4}{3}, \frac{7}{3}\right) \approx (1.33, 2.33)
                \end{equation*}
                
                \item Cluster 2 (P4, P5, P6):
                \begin{equation*}
                    \text{New centroid 2} = \left(\frac{5 + 6 + 6}{3}, \frac{4 + 5 + 7}{3}\right) = \left(\frac{17}{3}, \frac{16}{3}\right)\approx (5.67, 5.33)
                \end{equation*}
            \end{itemize}

            \item The assignment step in the second iteration is described in Table~\ref{tab:kmeans_iteration2}.
            \begin{table}[h!]
            \centering
            \begin{tabular}{|c|c|c|c|}
            \hline
            \textbf{Point} & \textbf{Distance to Centroid 1} & \textbf{Distance to Centroid 2} & \textbf{Assigned Cluster} \\ \hline
            P1 & $0.67$ & $8$ & Cluster 1 \\ \hline
            P2 & $2$ & $6$ & Cluster 1 \\ \hline
            P3 & $2$ & $8$ & Cluster 1 \\ \hline
            P4 & $5.33$ & $2$ & Cluster 2 \\ \hline
            P5 & $7.33$ & $0.67$ & Cluster 2 \\ \hline
            P6 & $9.33$ & $2$ & Cluster 2 \\ \hline
            \end{tabular}
            \caption{Assignment Step - Iteration 2}
            \label{tab:kmeans_iteration2}
            \end{table}

            \item Since the cluster assignments did not change from the first to the second iteration, the algorithm has converged. The final clusters are:
            \begin{itemize}
                \item Cluster 1: P1, P2, P3
                \item Cluster 2: P4, P5, P6
            \end{itemize}
        \end{enumerate}
    \end{ejercicio}

    \newpage

    \section*{Part 3: Kahn's Algorithm, Entropy \& Random Forest (22 pts.)}

    \begin{ejercicio}[Kahn's Algorithm (6 pts.)]
        Name and explain the steps of Kahn's Algorithm.

        Kahn's Algorithm is a method for performing topological sorting on a directed acyclic graph (DAG). The steps of Kahn's Algorithm are as follows:
        \begin{enumerate}
            \item Compute In-Degree: Calculate the in-degree (the number of incoming edges) for each vertex in the graph. This can be done by iterating through all edges and counting how many edges point to each vertex.
            \item Initialize Queue: Create a queue (or list) and enqueue all vertices with an in-degree of zero. These vertices have no dependencies and can be processed first.
            \item Process Vertices: While the queue is not empty, perform the following steps:
            \begin{enumerate}
                \item Dequeue a vertex from the queue and add it to the topological order list.
                \item For each outgoing edge from the dequeued vertex, reduce the in-degree of the target vertex by one. If the in-degree of the target vertex becomes zero, enqueue it into the queue.
            \end{enumerate}
            \item Check for Cycles: After processing all vertices, if the topological order list contains all the vertices in the graph, then a valid topological ordering has been found. If not, the graph contains a cycle, and topological sorting is not possible.
        \end{enumerate}
    \end{ejercicio}

    \begin{ejercicio}[Feature Types (2 pts.)]
        Consider Table~\ref{tab:customer_dataset2}. Which features are categorical and which are nominal?

        \begin{table}[h!]
        \centering
        \begin{tabular}{|l|l|l|l|}
        \hline
        \textbf{Client} & \textbf{Age} & \textbf{Ad} & \textbf{Buy} \\ \hline
        1 & young & yes & yes \\ \hline
        2 & young & no & no \\ \hline
        3 & medium & yes & yes \\ \hline
        4 & old & yes & yes \\ \hline
        5 & old & yes & no \\ \hline
        6 & medium & yes & yes \\ \hline
        \end{tabular}
        \caption{Customer Dataset}
        \label{tab:customer_dataset2}
        \end{table}

        In Table~\ref{tab:customer_dataset2}, all of the features are categorical.
        \begin{itemize}
            \item \textbf{Client}: It is nominal, as it represents unique identifiers for each client without any inherent order.
            \item \textbf{Age}: It is categorical and ordinal, as it represents age groups (young, medium, old) with a natural order.
            \item \textbf{Ad}: It is categorical and nominal, as it represents a binary choice (yes or no) without any inherent order.
            \item \textbf{Buy}: It is categorical and nominal, as it represents a binary choice (yes or no) without any inherent order.
        \end{itemize}
    \end{ejercicio}

    \begin{ejercicio}[Entropy Calculation (9 pts.)]
        Which feature should you split on first in Table~\ref{tab:customer_dataset2} if the column \textit{Buy} is the target? Calculate the entropy for each feature using the formula:
        \[ S(x) = -\sum_{i=1}^{N} p(x_i) \cdot \log_2(p(x_i)) \]

        Let's calculate the entropy for each feature in Table~\ref{tab:customer_dataset2} with respect to the target variable \textit{Buy}.
        \begin{itemize}
            \item Let's first consider spplit by \textbf{Age}.
            \begin{itemize}
                \item For \textit{Age = young}: 2 instances (1 yes, 1 no)
                \[ S(\text{young}) = -\left(\frac{1}{2} \log_2 \frac{1}{2} + \frac{1}{2} \log_2 \frac{1}{2}\right) = -\log_2 \frac{1}{2} = \log_2(2) = 1 \]
                \item For \textit{Age = medium}: 2 instances (2 yes, 0 no)
                
                Since it is the case of a pure node, the entropy is $S(\text{medium}) = 0$.
                \item For \textit{Age = old}: 2 instances (1 yes, 1 no)
                
                Since it is the same division as for \textit{young}, the entropy is $S(\text{old}) = 1$.
            \end{itemize}

            Therefore, the weighted entropy for \textit{Age} is:
            \[ S(\text{Age}) = \frac{2}{6} \cdot 1 + \frac{2}{6} \cdot 0 + \frac{2}{6} \cdot 1 = \frac{2}{6}\left(1 + 0 + 1\right) = \frac{2}{6} \cdot 2 = \frac{4}{6} = \frac{2}{3} \approx 0.667 \]

            \item Now let's consider split by \textbf{Ad}.
            \begin{itemize}
                \item For \textit{Ad = yes}: 5 instances (4 yes, 1 no)
                \[ S(\text{yes}) = -\left(\frac{4}{5} \log_2 \frac{4}{5} + \frac{1}{5} \log_2 \frac{1}{5}\right) = -\left(0.8 \cdot \log_2 0.8 + 0.2 \cdot \log_2 0.2\right) \approx 0.722 \]
                \item For \textit{Ad = no}: 1 instance (0 yes, 1 no)
                
                Since it is the case of a pure node, the entropy is $S(\text{no}) = 0$.
            \end{itemize}

            Therefore, the weighted entropy for \textit{Ad} is:
            \[ S(\text{Ad}) = \frac{5}{6} \cdot 0.722 + \frac{1}{6} \cdot 0 = 0.602 \]
        \end{itemize}

        Since the entropy for \textit{Ad} (0.602) is lower than that for \textit{Age} (0.667), we should split on the feature \textbf{Ad} first, as it provides a better separation of the target variable \textit{Buy}.
    \end{ejercicio}

    \begin{ejercicio}[Information Gain Range (2 pts.)]
        What is the range of values for Information Gain? Specify its minimum and maximum values.

        The formula for Information Gain (IG) is given by:
        \[ IG(T, A) = S(T) - \sum_{v \in \text{Values}(A)} \frac{|T_v|}{|T|} S(T_v) \]

        The range of values for Information Gain is as follows:
        \begin{itemize}
            \item Minimum Value: The minimum value of Information Gain is $0$. This occurs when the feature $A$ does not provide any information about the target variable $T$, meaning that the entropy of the target variable remains unchanged after the split.
            \item Maximum Value: The maximum value of Information Gain is equal to the entropy of the target variable $S(T)$. This occurs when the feature $A$ perfectly separates the data into pure subsets, resulting in a complete reduction of entropy.
        \end{itemize}
        Therefore, the range of Information Gain is $[0, S(T)]$, where $S(T)$ is the entropy of the target variable before the split.
    \end{ejercicio}

    \begin{ejercicio}[Random Forest Classification (3 pts.)]
        Classify the point $(F_1 = 1, F_2 = 4, F_3 = 7)$ using the Random Forest in Figure~\ref{fig:random_forest2}. Provide the explicit result for each tree as well as the final aggregated result. Draw the explicit decision path. Note: the left branch of a node represents \textbf{true}, and the right branch represents \textbf{false}.

        \begin{figure}[h!]
            \centering
            \resizebox{\textwidth}{!}{%
            \begin{tikzpicture}[
                node/.style={circle, draw=black, minimum size=1.2cm, inner sep=0pt, font=\scriptsize\bfseries, align=center},
                leaf/.style={circle, draw=black, minimum size=0.85cm, inner sep=0pt, font=\scriptsize\bfseries},
                edge from parent/.style={draw, -Stealth, ultra thick},
                tree label/.style={font=\small\bfseries, yshift=0.35cm}
            ]

            % ==================== TREE 1 ====================
            \begin{scope}[shift={(0,0)}]
                \node[tree label] at (0,0) {1};
                \node[node] at (0,-0.6) {$F_2 \leq 4.5$}
                    [sibling distance=3.2cm, level distance=1.4cm]
                    child { node[node] {$F_1 \leq 1$}
                        [sibling distance=1.6cm, level distance=1.4cm]
                        child { node[node] {$F_3 > 6$}
                            [sibling distance=1.2cm, level distance=1.4cm]
                            child { node[leaf] {C1} }
                            child { node[leaf] {C3} }
                        }
                        child { node[leaf] {C4} }
                    }
                    child { node[node] {$F_1 > 1$}
                        [sibling distance=1.6cm, level distance=1.4cm]
                        child { node[node] {$F_1 < -1$}
                            [sibling distance=1.2cm, level distance=1.4cm]
                            child { node[leaf] {C3} }
                            child { node[leaf] {C2} }
                        }
                        child { node[leaf] {C1} }
                    };
            \end{scope}

            % ==================== TREE 2 ====================
            \begin{scope}[shift={(6.2,0)}]
                \node[tree label] at (0,0) {2};
                \node[node] at (0,-0.6) {$F_2 > 2$}
                    [sibling distance=2.6cm, level distance=1.4cm]
                    child { node[node] {$F_1 < 0$}
                        [sibling distance=1.5cm, level distance=1.4cm]
                        child { node[node] {$F_3 < 9$}
                            [sibling distance=1.2cm, level distance=1.4cm]
                            child { node[leaf] {C1} }
                            child { node[leaf] {C2} }
                        }
                        child { node[leaf] {C1} }
                    }
                    child { node[node] {$F_2 \leq 1$}
                        [sibling distance=1.2cm, level distance=1.4cm]
                        child { node[leaf] {C4} }
                        child { node[leaf] {C2} }
                    };
            \end{scope}

            % ==================== TREE 3 ====================
            \begin{scope}[shift={(11.8,0)}]
                \node[tree label] at (0,0) {3};
                \node[node] at (0,-0.6) {$F_3 \leq 10$}
                    [sibling distance=1.8cm, level distance=1.4cm]
                    child { node[node] {$F_3 \leq 5$}
                        [sibling distance=1.3cm, level distance=1.4cm]
                        child { node[leaf] {C3} }
                        child { node[leaf] {C1} }
                    }
                    child { node[leaf] {C2} };
            \end{scope}

            % ==================== TREE 4 ====================
            \begin{scope}[shift={(7.2,-5.2)}]
                \node[tree label] at (0,0) {4};
                \node[node] at (0,-0.6) {$F_1 \geq 2$}
                    [sibling distance=1.4cm, level distance=1.4cm]
                    child { node[leaf] {C2} }
                    child { node[leaf] {C5} };
            \end{scope}

            % ==================== TREE 5 ====================
            \begin{scope}[shift={(10.8,-5.2)}]
                \node[tree label] at (0,0) {5};
                \node[node] at (0,-0.6) {$F_2 \geq 4$}
                    [sibling distance=1.4cm, level distance=1.4cm]
                    child { node[leaf] {C4} }
                    child { node[leaf] {C1} };
            \end{scope}

            \end{tikzpicture}
            }
            \caption{Random Forest Classifier}
            \label{fig:random_forest2}
        \end{figure}

        Let's classify the point $(F_1 = 1, F_2 = 4, F_3 = 7)$ using the Random Forest in Figure~\ref{fig:random_forest2}.
        \begin{itemize}
            \item According to Tree 1, the point is classified as C1.
            \item According to Tree 2, the point is classified as C1.
            \item According to Tree 3, the point is classified as C1.
            \item According to Tree 4, the point is classified as C5.
            \item According to Tree 5, the point is classified as C4.
        \end{itemize}

        The final aggregated result is determined by majority voting among the trees. Therefore, the final classification for the point $(F_1 = 1, F_2 = 4, F_3 = 7)$ is C1, as it received the most votes from the individual trees.
    \end{ejercicio}

    \newpage

    \section*{Part 4: NLP \& Forward Propagation (8 pts.)}

    \begin{ejercicio}[Bag of Words Encoding (4 pts.)]
        Consider the sentence \textit{``Nobody makes it better than me, nobody.''} and the vocabulary
        $$\{A, be, better, easy, is, it, make, makes, me, nobody, than, then\}$$
        
        Encode the sentence using frequency-based Bag of Words.\\

        The frequency of each word in the sentence \textit{``Nobody makes it better than me, nobody.''} based on the given vocabulary is as described in Figure~\ref{fig:bow_encoding}. The encoded representation is a vector where each element corresponds to the frequency of the respective word in the vocabulary.
        \begin{table}[h!]
        \centering
        \begin{tabular}{|c|c|c|c|c|c|c|c|c|c|c|c|}
        \hline
        \textbf{A} & \textbf{be} & \textbf{better} & \textbf{easy} & \textbf{is} & \textbf{it} & \textbf{make} & \textbf{makes} & \textbf{me} & \textbf{nobody} & \textbf{than} & \textbf{then} \\ \hline
        0 & 0 & 1 & 0 & 0 & 1 & 0 & 1 & 1 & 2 & 1 & 0 \\ \hline
        \end{tabular}
        \caption{Frequency-based Bag of Words Encoding}
        \label{fig:bow_encoding}
        \end{table}
    \end{ejercicio}

    \begin{ejercicio}[Forward Propagation Computation (4 pts.)]
        Given the small neural network below, write the calculated output values (after applying the activation functions) into the corresponding boxes. The weights $w$ are provided for all connections.

        \begin{figure}[h!]
            \centering
            \resizebox{\textwidth}{!}{%
            \begin{tikzpicture}[
                neuron/.style={circle, draw, fill=black, text=white, minimum size=1.3cm, font=\small\sffamily\bfseries, align=center},
                box/.style={rectangle, draw, minimum size=0.65cm, thick, fill=white},
                edge/.style={-Stealth, thick},
                weight/.style={font=\small, fill=white, inner sep=1.5pt, rounded corners=1pt, sloped}
            ]

            % --- Nodes Placement ---
            % Input
            \node[font=\large\bfseries] (in) at (-1.5,0) {7};

            % Layer 1: Heaviside
            \node[neuron] (H) at (1,0) {Heaviside};
            \node[box] (H_box) at (2.4,0) {1};

            % Layer 2: ReLUs
            \node[neuron] (R1) at (6, 2.8) {ReLU};
            \node[box] (R1_box) at (6, 3.9) {1};

            \node[neuron] (R2) at (6, 0) {ReLU};
            \node[box] (R2_box) at (6, 1.1) {0};

            \node[neuron] (R3) at (6, -2.8) {ReLU};
            \node[box] (R3_box) at (6, -3.9) {1};

            % Layer 3: Sign
            \node[neuron] (S1) at (11.5, 1.6) {Sign};
            \node[box] (S1_box) at (11.5, 2.7) {1};

            \node[neuron] (S2) at (11.5, -1.6) {Sign};
            \node[box] (S2_box) at (11.5, -2.7) {1};

            % Layer 4: Output ReLU
            \node[neuron] (Out) at (16, 0) {ReLU};
            \node[box] (Out_box) at (17.4, 0) {0};


            % --- Connections & Weights ---

            % Input -> Heaviside
            \draw[edge] (in) -- (H);

            % Heaviside Box -> ReLUs
            \draw[edge] (H_box) -- node[weight, above] {$w_2 = 1$} (R1);
            \draw[edge] (H_box) -- node[weight, above] {$w_1 = -2$} (R2);
            \draw[edge] (H_box) -- node[weight, below] {$w_3 = 1$} (R3);

            % ReLUs -> Sign Layer (Weights placed clearly near start or near end, sloped along the line)
            \draw[edge] (R1) -- node[weight, pos=0.2, above] {$w_1 = 2$} (S1);
            \draw[edge] (R1) -- node[weight, pos=0.2, below] {$w_2 = 2$} (S2);

            \draw[edge] (R2) -- node[weight, pos=0.7, above] {$w_3 = 2$} (S1);
            \draw[edge] (R2) -- node[weight, pos=0.7, below] {$w_4 = 2$} (S2);

            \draw[edge] (R3) -- node[weight, pos=0.2, above] {$w_5 = 2$} (S1);
            \draw[edge] (R3) -- node[weight, pos=0.2, below] {$w_6 = 2$} (S2);

            % Sign Layer -> Output
            \draw[edge] (S1) -- node[weight, above] {$w_1 = -2$} (Out);
            \draw[edge] (S2) -- node[weight, below] {$w_2 = 2$} (Out);

            \end{tikzpicture}
            }
            \caption{Forward Propagation Network Schema}
        \end{figure}
    \end{ejercicio}

\end{document}